\makeatletter \def\input@path{{../../}} \makeatother
\documentclass[11pt]{article}

\usepackage[utf8]{inputenc}
\usepackage{graphicx}
\usepackage[dvipsnames]{xcolor}
\usepackage{color}
\usepackage[margin=1in]{geometry}
\usepackage{hyperref}
\usepackage{tikz}
\usepackage{amsmath}
\usepackage{commath}
\usepackage{relsize}
\usepackage{centernot}
\usepackage{amssymb}
\usepackage{amsfonts}
\usepackage{graphicx}
\usepackage{textcomp}
\usepackage{gensymb}
\usepackage{enumitem}
\usepackage{siunitx}
\usepackage{fancyhdr}
\usepackage{floatrow}
\usepackage{wrapfig}
\usepackage[labelfont=bf]{caption}
\usepackage{mathtools}
\usepackage{pgfplots}
\usepackage{multicol}
\usepackage[misc]{ifsym}
\usepackage{tabularx}
\usepackage{empheq}
\usepackage{tkz-tab}
\usepackage{xpatch}
\usepackage{caption}
\usepackage{subcaption}
\usepackage{tabularray}
\usepackage{esint}
\usepackage{amsthm}
\usepackage{thmtools}
\RequirePackage[most]{tcolorbox}

\swapnumbers

\makeatletter
\declaretheoremstyle[
  headfont= \sffamily\bfseries\color{cyan!50!black},
  headpunct={\sffamily\bfseries.},
  postheadspace=0.5em,
  notefont=\sffamily\bfseries,
  headformat=\NAME\NUMBER\let\thmt@space\@empty\NOTE,
  bodyfont=\mdseries\slshape,
  spaceabove=12pt,spacebelow=12pt,
]{thmstyle}
\makeatother

\theoremstyle{thmstyle}
\declaretheorem[name=Definici\'on]{theorem}
\newtheorem{definition}[theorem]{Definici\'on}

\definecolor{lightcyan}{rgb}{0.88, 1.0, 1.0}
\colorlet{mythmback}{lightcyan!40!white}

\tcolorboxenvironment{theorem}{
  colback=mythmback,coltitle=blue,colframe=mythmback,
  left=0pt,right=0pt,
  top=0pt,bottom=0pt,
  before skip=10pt, after skip=10pt,
}

\xpatchcmd{\tkzTabLine}{$0$}{$\bullet$}{}{}
\tikzset{t style/.style={style=solid}}
\pgfplotsset{compat=newest}

\newtheorem{manualtheoreminner}{}
\newenvironment{manualtheorem}[1]{%
  \renewcommand\themanualtheoreminner{#1}%
  \manualtheoreminner
}{\endmanualtheoreminner}

\renewcommand{\pd}[2]{\frac{\partial #1}{ \partial #2}}
\newcommand{\td}[2]{\frac{\mathrm{d} #1}{ \mathrm{d} #2}}
\newcommand{\pdd}[2]{\frac{\partial^2 #1}{ \partial #2 ^2}}
\newcommand{\pddc}[3]{\frac{\partial^2 #1}{ \partial #2 \partial #3}}
\newcommand{\solution}[0]{\large{\textbf{Soluci\'on}}\normalsize}

\def\sca   #1{\mbox{\rm{#1}}{}}
\def\mat   #1{\mbox{\boldmath $\mathsf #1$}}
\def\vec   #1{\mbox{\boldmath $#1$}{}}
\def\ten   #1{\mbox{\boldmath $#1$}{}}
\def\ltr   #1{\mbox{\sf{#1}}}
\def\bltr  #1{\mbox{\sffamily{\bfseries{{#1}}}}}

\DeclareMathOperator{\dive}{div}
\DeclareMathOperator{\trace}{trace}
\DeclareMathOperator{\tr}{tr}
\DeclareMathOperator{\symm}{symm}
\DeclareMathOperator{\sk}{skew}
\DeclareMathOperator{\grad}{grad}
\DeclareMathOperator{\Grad}{Grad}
\DeclareMathOperator{\curl}{curl}
\DeclareMathOperator{\Curl}{Curl}
\DeclareMathOperator*{\argmin}{\arg\!\min}

\def\dx{\mbox{\(\,\mathrm{d}x\)}}
\def\ds{\mbox{\(\,\mathrm{d}s\)}}
\def\dS{\mbox{\(\,\mathrm{d}S\)}}
\def\dV{\mbox{\(\,\mathrm{d}V\)}}
\def\dv{\mbox{\(\,\mathrm{d}v\)}}
\def\R{\mathbb R}
\def\C{\mathbb C}
\def\N{\mathbb N}
\def\Q{\mathbb Q}
\def\Z{\mathbb Z}
\def\D{\mathcal D}

\usepackage{mdframed}
\mdfdefinestyle{MyFrame}{%
    linecolor=black,
    linewidth=1pt,
    outerlinewidth=1pt,
    roundcorner=20pt,
    innertopmargin=\baselineskip,
    innerbottommargin=\baselineskip,
    innerrightmargin=20pt,
    innerleftmargin=20pt,
    splittopskip=2\baselineskip,
    backgroundcolor=gray!50!white}

\setlength\textwidth{7in}
\setlength\parindent{0pt}
\oddsidemargin=-0.5cm
\pagestyle{fancy}
\newcommand*{\vertbar}{\rule[-1ex]{0.5pt}{2.5ex}}
\setlength{\headheight}{13.59999pt}

\renewcommand{\headrulewidth}{0.1pt}
\renewcommand{\footrulewidth}{0.1pt}
\renewcommand{\figurename}{Figura}
\renewcommand\refname{Referencias}
\renewcommand{\thesection}{\arabic{section}.}
\renewcommand{\thesubsection}{\thesection\arabic{subsection}.}
\renewcommand{\labelitemi}{{\raisebox{.3\height}{\scalebox{0.6}{$\blacklozenge$}}}}
\renewcommand*\contentsname{\'Indice}

\newcommand{\p}{\vspace{10pt}}

\AtBeginDocument{\vspace*{-3.2\baselineskip}}
\textheight=665pt

\lhead{\scshape Pontificia Universidad Cat\'olica de Chile}
\rhead{\scshape IMC UC}
\cfoot{\thepage}

\usepackage{footnote}
\usepackage{etoolbox}
\newtoggle{indefintion}
\togglefalse{indefintion}
\pretocmd{\footnote}{\iftoggle{indefintion}{\stepcounter{footnote}}{\relax}}{}{}


\begin{document}
    \begin{flushleft}
        \small
        \vspace{1pt}
        \textbf{IMT3130} \hfill
        \textbf{09/05/2026}
    \end{flushleft}
    \begin{center}
        \LARGE\textbf{Ayudant\'ia 7}\\
        \vspace{3pt}
        \normalsize\textbf{Interpolaci\'on, elementos finitos compatibles e inf--sup}\\
        \normalsize Ayudante: Basti\'an Herrera \Letter{\hspace{1pt}: \texttt{biherrera@uc.cl}}
        \rule{\textwidth}{0.05mm}
    \end{center}

\section*{Problemas}

\begin{enumerate}[label=\textbf{\arabic*.}, leftmargin=*]

    \item \textbf{C\'ea, interpolaci\'on y Aubin--Nitsche para reacci\'on--difusi\'on.} Sea $\eta>0$ fijo y considere
    \begin{equation*}
    \begin{aligned}
        -u''+\eta u &= f &&\text{en }(0,1),\\
        u(0)&=u(1)=0,
    \end{aligned}
    \end{equation*}
    con $f\in L^2(0,1)$. La formulaci\'on d\'ebil es: hallar $u\in V:=H^1_0(0,1)$ tal que
    \begin{equation*}
        a_\eta(u,v)=F(v)\qquad \forall v\in V,
        \qquad
        a_\eta(w,v)=\int_0^1 w'v'\,dx+\eta\int_0^1wv\,dx,
        \qquad
        F(v)=\int_0^1fv\,dx.
    \end{equation*}
    Defina la norma de energ\'ia $\|v\|_{a_\eta}^2:=a_\eta(v,v)$. Sea $\mathcal T_h$ una malla uniforme y, para $r=1,2$, sea
    \begin{equation*}
        V_h^{(r)}=\{v_h\in C^0([0,1]):\ v_h|_K\in\mathbb P_r(K)\ \forall K\in\mathcal T_h,\ v_h(0)=v_h(1)=0\}.
    \end{equation*}
    Sea $u_h^{(r)}\in V_h^{(r)}$ la soluci\'on de Galerkin. Suponga que la soluci\'on exacta satisface $u\in H^3(0,1)$ y use las estimaciones de interpolaci\'on nodal
    \begin{equation*}
        |u-\Pi_h^{(r)}u|_{H^1(0,1)}\le Ch^r|u|_{H^{r+1}(0,1)},
        \qquad
        \|u-\Pi_h^{(r)}u\|_{L^2(0,1)}\le Ch^{r+1}|u|_{H^{r+1}(0,1)}.
    \end{equation*}
    \begin{enumerate}[label=(\alph*)]
        \item Pruebe la ortogonalidad de Galerkin y deduzca la estimaci\'on de C\'ea en norma de energ\'ia.
        \item Compare los errores de energ\'ia esperados para $P_1$ y $P_2$.
        \item Use el argumento de Aubin--Nitsche para obtener estimaciones en norma $L^2$ para $P_1$ y $P_2$. Explique por qu\'e el resultado en $L^2$ es estrictamente mejor que el resultado en energ\'ia.
    \end{enumerate}

    \item \textbf{Dos triples de Ciarlet: unisolvencia de $P_2$ y $RT_0$.} Recuerde que un elemento finito de Ciarlet es un triple $(K,\mathcal P,\Sigma)$, donde $K$ es el dominio local, $\mathcal P$ es el espacio local y $\Sigma$ es una familia de grados de libertad que debe determinar un\'ivocamente a cada funci\'on de $\mathcal P$.
    \begin{enumerate}[label=(\alph*)]
        \item En el intervalo de referencia $\widehat K=[-1,1]$, considere $\mathcal P=\mathbb P_2(\widehat K)$ y
        \begin{equation*}
            \ell_1(p)=p(-1),\qquad \ell_2(p)=p(0),\qquad \ell_3(p)=p(1).
        \end{equation*}
        Pruebe la unisolvencia y construya las tres funciones de base locales.

        \item En el tri\'angulo de referencia
        \begin{equation*}
            \widehat K=\operatorname{conv}\{(0,0),(1,0),(0,1)\},
        \end{equation*}
        considere
        \begin{equation*}
            RT_0(\widehat K)=\{\vec q(x,y)=(\alpha,\beta)+\gamma(x,y):\ \alpha,\beta,\gamma\in\mathbb R\}.
        \end{equation*}
        Los grados de libertad son los flujos normales
        \begin{equation*}
            \ell_i(\vec q)=\int_{e_i}\vec q\cdot\vec n_i\,ds,
            \qquad i=1,2,3,
        \end{equation*}
        donde $e_i$ son las tres aristas y $\vec n_i$ las normales exteriores. Pruebe la unisolvencia en este tri\'angulo.

        \item Explique qu\'e continuidad global induce el ensamblaje de $P_2$ y qu\'e continuidad global induce el ensamblaje de $RT_0$.
    \end{enumerate}

    \item \textbf{Inf--sup continuo versus inf--sup discreto.} Considere el problema de transporte--reacci\'on
    \begin{equation*}
        u'+u=f\quad\text{en }(0,1),
        \qquad
        u(0)=0.
    \end{equation*}
    Sean
    \begin{equation*}
        U=\{u\in H^1(0,1):u(0)=0\},
        \qquad
        V=L^2(0,1),
        \qquad
        a(u,v)=\int_0^1(u'+u)v\,dx.
    \end{equation*}
    \begin{enumerate}[label=(\alph*)]
        \item Muestre que $a$ no es coerciva con respecto a la norma de $H^1(0,1)$.
        \item Pruebe que el problema continuo satisface una condici\'on inf--sup en $U\times V$.
        \item Considere la elecci\'on discreta deliberadamente incompatible
        \begin{equation*}
            U_h=\operatorname{span}\{x,x^2\},
            \qquad
            V_h=\operatorname{span}\{1\}.
        \end{equation*}
        Pruebe que la condici\'on inf--sup discreta falla. Interprete el resultado en relaci\'on con el an\'alisis de Galerkin para problemas no coercivos.
    \end{enumerate}

\end{enumerate}

\end{document}
